%; whizzy-master slides.tex
%; whizzy frame -pdf ./skimopen 

\usepackage{xspace} % To get the right spacings in front of : and so on
\usepackage[english]{babel}
\usepackage{subfigure}
\usepackage{bussproofs}
\usepackage{pgf}
\usepackage{colortbl}
\usepackage{xcolor}
\usepackage{abbrevs}
\usepackage[color]{coqdoc}
\def\coqlibrary#1#2{}

\usepackage{hyperref}
\usepackage{coq}
\usepackage{code}
\usepackage{codecolor}
\usepackage{cond}
\usepackage{me}
\usepackage{tikz}
\usetikzlibrary{shapes}
\usepackage{pifont}
\usepackage{proof}
\usepackage{utf}
\usepackage{natbib}
\usepackage[OT1]{fontenc}
\usepackage{alltt}
\usepackage{amsthm}


%%% MMCG symbols

\newcommand{\idot}[1]{\mbox{$\bullet_{#1}$}}
\newcommand{\baredot}{\mbox{$\bullet$}}
\newcommand{\iback}[1]{\mbox{$\backslash_{#1}$}}
\newcommand{\back}{\mbox{$\backslash$}}
\newcommand{\islash}[1]{\mbox{$/_{#1}$}}
\newcommand{\icomma}[3]{\mbox{$(#1\,,\,#2)_{#3}$}}
%\newcommand{\icomma}[3]{\mbox{$#1\,\mathbin{\circ}_{#3}\,#2$}}
%\newcommand{\comma}{\mbox{$\mathbin{\circ}$}}
\newcommand{\comma}{\mbox{,}}

\newcommand{\slashe}{\mbox{\textbf{$/_{\textbf{E}}$}}}
\newcommand{\slashi}{\mbox{\textbf{$/_{\textbf{I}}$}}}
\newcommand{\backe}{\mbox{\textbf{$\backslash_{\textbf{E}}$}}}
\newcommand{\backi}{\mbox{\textbf{$\backslash_{\textbf{I}}$}}}

\newcommand{\dote}{\mbox{\textbf{$\baredot_{\textbf{E}}$}}}
\newcommand{\doti}{\mbox{\textbf{$\baredot_{\textbf{I}}$}}}
\newcommand{\diame}{\mbox{\textbf{$\Diamond_{\textbf{E}}$}}}
\newcommand{\diami}{\mbox{\textbf{$\Diamond_{\textbf{I}}$}}}
\newcommand{\struct}{\mbox{\textbf{struct}}}
\newcommand{\axrule}{\mbox{\textbf{Ax}}}
\newcommand{\boxe}{\mbox{\textbf{$\Box_{\textbf{E}}$}}}
\newcommand{\boxi}{\mbox{\textbf{$\Box_{\textbf{I}}$}}}
\newcommand{\AssD}{\textbf{L$_\Diamond$}}
\newcommand{\ComD}{\textbf{P$_\Diamond$}}

\newcommand{\marginok}[1]{\marginpar{\raggedright OK:#1}}
\newcommand{\tab}{{\null\hskip1cm}}
%\newcommand{\Ltac}{\mbox{\emph{$\cal L$}tac}}
\newcommand{\coq}{\mbox{\emph{Coq}}}
\newcommand{\pcoq}{\mbox{\emph{Pcoq}}}
\newcommand{\grail}{\mbox{\emph{Grail}}}
\newcommand{\lcf}{\mbox{\emph{LCF}}}
\newcommand{\hol}{\mbox{\emph{HOL}}}
\newcommand{\pvs}{\mbox{\emph{PVS}}}
\newcommand{\icharate}{\mbox{\emph{Icharate}}}
\newcommand{\isabelle}{\mbox{\emph{Isabelle}}}
\newcommand{\prolog}{\mbox{\emph{Prolog}}}
\newcommand{\goalbar}{\tt{}{\color{black}------------------------------------}\it}
\newcommand{\gallina}{\mbox{\emph{Gallina}}}
\newcommand{\joker}{\texttt{\_}}
\newcommand{\eprime}{\(\e^{\prime}\)}
\newcommand{\Ztype}{\citecoq{Z}}
\newcommand{\propsort}{\citecoq{Prop}}
\newcommand{\setsort}{\citecoq{Set}}
\newcommand{\typesort}{\citecoq{Type}}
\newcommand{\ocaml}{\mbox{\emph{OCAML}}}
\newcommand{\haskell}{\mbox{\emph{Haskell}}}
\newcommand{\why}{\mbox{\emph{Why}}}
\newcommand{\Pascal}{\mbox{\emph{Pascal}}}

\newcommand{\ml}{\mbox{\emph{ML}}}

\newcommand{\scheme}{\mbox{\emph{Scheme}}}
\newcommand{\lisp}{\mbox{\emph{Lisp}}}

\newcommand{\implarrow}{\mbox{$\Rightarrow$}}
\newcommand{\metavar}[1]{?#1}
\newcommand{\notincoq}[1]{#1}
\newcommand{\coqscope}[1]{\%#1}
\newcommand{\arrow}{\mbox{$\rightarrow$}}
\newcommand{\fleche}{\arrow}
\newcommand{\funarrow}{\mbox{$\Rightarrow$}}
\newcommand{\ltacarrow}{\funarrow}
% \newcommand{\coqand}{\mbox{\(\wedge\)}}
% \newcommand{\coqor}{\mbox{\(\vee\)}}
% \newcommand{\coqnot}{\mbox{\(\neg\)}}
\newcommand{\coqand}{\texttt{/\textbackslash}}
%\newcommand{\coqor}{\texttt{\textbackslash/}}
\newcommand{\coqnot}{\mbox{\(\neg\)}}
\newcommand{\hide}[1]{}
\newcommand{\hidedots}[1]{...}
\newcommand{\sig}[3]{\texttt{\{}#1\texttt{:}#2 \texttt{|} #3\texttt{\}}}
\renewcommand{\neg}{\mbox{$\sim$}}

%%% Operateurs, etc.
\newcommand{\impl}{\mbox{$\rightarrow$}}
\newcommand{\appli}[2]{\mbox{\tt{#1 #2}}}
\newcommand{\applis}[1]{\mbox{\texttt{#1}}}
\newcommand{\abst}[3]{\mbox{\tt{fun #1:#2 \funarrow #3}}}
\newcommand{\coqle}{\mbox{$\leq$}}
\newcommand{\coqge}{\mbox{$\geq$}}
\newcommand{\coqdiff}{\mbox{$\neq$}}
\newcommand{\coqiff}{\mbox{$\leftrightarrow$}}
\newcommand{\prodsym}{\mbox{\(\forall\,\)}}
\newcommand{\exsym}{\mbox{\(\exists\,\)}}
\newcommand{\included}{\mbox{\(\subseteq\)}}

\newcommand{\substsign}{/}
\newcommand{\subst}[3]{\mbox{#1\{#2\substsign{}#3\}}}
\newcommand{\anoabst}[2]{\mbox{\tt[#1]#2}}
\newcommand{\letin}[3]{\mbox{\tt let #1:=#2 in #3}}
\newcommand{\prodep}[3]{\mbox{\tt \(\forall\,\)#1:#2,$\,$#3}}
\newcommand{\prodplus}[2]{\mbox{\tt\(\forall\,\)$\,$#1,$\,$#2}}
\renewcommand{\dom}[1]{\textrm{dom}(#1)} % domaine d'un contexte (log function)
\newcommand{\norm}[1]{\textrm{n}(#1)} % forme normale (log function)
\newcommand{\coqZ}[1]{\mbox{\tt{`#1`}}}
\newcommand{\coqnat}[1]{\mbox{\tt{#1}}}
\newcommand{\coqcart}[2]{\mbox{\tt{#1*#2}}}
\newcommand{\alphacong}{\mbox{$\,\cong_{\alpha}\,$}} % alpha-congruence
\newcommand{\betareduc}{\mbox{$\,\rightsquigarrow_{\!\beta}$}\,} % beta reduction
%\newcommand{\betastar}{\mbox{$\,\Rightarrow_{\!\beta}^{*}\,$}} % beta reduction
\newcommand{\deltareduc}{\mbox{$\,\rightsquigarrow_{\!\delta}$}\,} % delta reduction
\newcommand{\dbreduc}{\mbox{$\,\rightsquigarrow_{\!\delta\beta}$}\,} % delta,beta reduction
\newcommand{\ireduc}{\mbox{$\,\rightsquigarrow_{\!\iota}$}\,} % delta,beta reduction

\newcommand{\slam}[2] {\mbox{$\lambda_#1\;#2$}}
\newcommand{\srho}[2] {\mbox{$\rho_#1\;#2$}}
\newcommand{\Vforall}[2] {\mbox{$\forall\,#1\,\in\,V(#2)$}}
\newcommand{\Eforall}[2] {\mbox{$\forall\,#1\,\in\,E(#2)$}}
\newcommand{\Vforone}[2] {\mbox{$\exists!\,#1\,\in\,V(#2)$}}
\newcommand{\Eforone}[2] {\mbox{$\exists!\,#1\,\in\,E(#2)$}}
\newcommand{\Vforsome}[2] {\mbox{$\exists\,#1\,\in\,V(#2)$}}
\newcommand{\Eforsome}[2] {\mbox{$\exists\,#1\,\in\,E(#2)$}}


% jugement de typage
\newcommand{\these}{\mbox{$\boldsymbol{\large \vdash}$}}
\newcommand{\rthese}[1]{\mbox{$\boldsymbol{\vdash_{#1}}$}}
\newcommand{\replace}[2]{\mbox{$#1[#2]$}}
\newcommand{\msubst}[2]{\mbox{$#1\{#2\}$}}
\newcommand{\disj}{\mbox{$\backslash/$}}
\newcommand{\conj}{\mbox{$/\backslash$}}
\newcommand{\deriv}[2]{\mbox{$#1\;\these\;#2$}}
\newcommand{\smalljuge}[3]{\mbox{$#1 \these #2 \boldsymbol{:} #3 $}}
%\newcommand{\juge}[3]{\mbox{$#1 \these #2 \boldsymbol{:} #3 $}}
\newcommand{\goal}[3]{\mbox{$#1,#2 \these^{\!\!\!?} #3  $}}
\newcommand{\sgoal}[2]{\mbox{$#1\,\these^{\!\!\!\!?}\, #2 $}}
\newcommand{\sequent}[2]{\mbox{$#1\;\these\; #2 $}}

\newcommand{\reduc}[5]{\mbox{$#1,#2 \these #3 \rhd_{#4}#5 $}}
\newcommand{\convert}[5]{\mbox{$#1,#2 \these #3 =_{#4}#5 $}}
\newcommand{\convorder}[5]{\mbox{$#1,#2 \these #3\leq _{#4}#5 $}}
\newcommand{\wouff}[2]{\mbox{$\emph{WF}(#1)[#2]$}}




\newcommand{\type}{\boldsymbol{:}}

% jugement absolu

%\newcommand{\ajuge}[2]{\mbox{$ \boldsymbol{\vdash} #1 : #2 $}}
\newcommand{\ajuge}[2]{\mbox{$\these #1 \boldsymbol{:} #2 $}}

%%% logique minimale
\newcommand{\propzero}{\mbox{$P_0$}} % types de Fzero

%%% logique propositionnelle classique
\newcommand {\ff}{\boldsymbol{f}} % faux
\newcommand {\vv}{\boldsymbol{t}} % vrai

\newcommand{\verite}{\mbox{$\cal{B}$}} % {\ff,\vv}
\newcommand{\sequ}[2]{\mbox{$#1 \vdash #2 $}} % sequent
\newcommand{\strip}[1]{#1^o} % enlever les variables d'un contexte



%%% tactiques
\newcommand{\decomp}{\delta} % decomposition
\newcommand{\recomp}{\rho} % recomposition

%%% divers
\newcommand{\cqfd}{\mbox{\textbf{cqfd}}}
\newcommand{\fail}{\mbox{\textbf{F}}}
\newcommand{\succes}{\mbox{$\blacksquare$}}
%%% Environnements


%% Fzero
\newcommand{\con}{\mbox{$\cal C$}}
\newcommand{\var}{\mbox{$\cal V$}}

\newcommand{\atomzero}{\mbox{${\cal A}_0$}} % types de base de Fzero
\newcommand{\typezero}{\mbox{${\cal T}_0$}} % types de Fzero
\newcommand{\termzero}{\mbox{$\Lambda_0$}} % termes de Fzero 
\newcommand{\conzero}{\mbox{$\cal C_0$}} % contextes de Fzero 

\newcommand{\buts}{\mbox{$\cal B$}} % buts

%%% for drawing terms
% abstraction [x:t]e
\newcommand{\PicAbst}[3]{\begin{bundle}{\bf abst}\chunk{#1}\chunk{#2}\chunk{#3}%
                        \end{bundle}}

% the same in DeBruijn form
\newcommand{\PicDbj}[2]{\begin{bundle}{\bf abst}\chunk{#1}\chunk{#2}
                       \end{bundle}}


% applications
\newcommand{\PicAppl}[2]{\begin{bundle}{\bf appl}\chunk{#1}\chunk{#2}%
                         \end{bundle}}

% variables
\newcommand{\PicVar}[1]{\begin{bundle}{\bf var}\chunk{#1}
                          \end{bundle}}

% constantes
\newcommand{\PicCon}[1]{\begin{bundle}{\bf const}\chunk{#1}\end{bundle}}

% arrows
\newcommand{\PicImpl}[2]{\begin{bundle}{\impl}\chunk{#1}\chunk{#2}%
                         \end{bundle}}



%%%% scripts coq
\newcommand{\prompt}{\mbox{\sl Coq $<\;$}}
\newcommand{\natquicksort}{\texttt{nat\_quicksort}}
\newcommand{\citecoq}[1]{\mbox{\texttt{#1}}}
\newcommand{\safeit}{\it}
\newcommand{\byrule}[1]{\mbox{$\underset{\scriptstyle #1}{\longrightarrow}$}}
\newcommand{\byrules}[1]{\mbox{$\underset{\scriptstyle #1}{\longrightarrow^*}$}}


\newcommand{\backback}{\mbox{$\backslash\backslash$}}

\renewcommand{\cons}{\mbox{::}}
\input{univmacros}

\def\vec#1{\overrightarrow{#1}}
\def\thetitle{Universe Polymorphism: Subtyping and Unification}
\def\thetextitleabbr{\thetitle}
\def\thetextitle{\thetextitleabbr}
\def\thesubject{Theoretical Computer Science}
\def\thekeywords{Coq, Universes}
\def\theevent{TYPEX Days\\
December 17th 2013\\
Paris}
\def\theeventabbr{}
\def\theauthorabbr{M. Sozeau}
\usepackage{prelude-beamer}

\author[\myabbrname{}]{{\sc \myname{}}}
\def\PI{\name{PI}}
\def\shaded#1{{\color{black!50}#1}}
 
\def\seq{\fCenter}

\def\lang{en}
\input{mathenv}
\setboolean{displayLabels}{true}

\pgfdeclareimage[height=1.4cm]{inria-logo}{inria-logo}

\pgfdeclareimage[height=1.4cm]{coq-biglogo}{barron-logo-pdf}
\pgfdeclareimage[height=0.4cm]{coq-smalllogo}{barron-logo}
\pgfdeclareimage[height=1cm,interpolate=true]{coq-logo}{barron-logo-pdf}

\pgfdeclareimage[height=7cm]{coq-share}{coq-share}

\def\imgheight{7.5cm}

\def\typea{\vdash}
\def\suba{\rightslice}
\def\indent{\coqdocindent}
\def\var{\coqdocvar}

\def\coqdocid#1{\coqdocvar{#1}}

%\renewcommand{\Type}[1]{\kw{Type}_{#1}}

\newcommand{\elt}[4]{\ensuremath{\constr{exist}_{#1,#2}~#3~#4}}
\newcommand{\eltlight}[2]{\ensuremath{\constr{exist}~#1~#2}}
\def\class{\kw{class}~}
\def\instance{\kw{instance}~}
\def\tclass#1{\coqdocind{#1}}
\def\where{\kw{where}~}
\def\module#1{\texttt{#1}}

\def\vector#1{\ind{vector}~#1}
\def\vnil{\constr{vnil}}
\def\vcons{\constr{vcons}}

\renewcommand{\bar}[1]{{\overline{#1}}}

\setlength{\coqdocbaseindent}{0em}

\renewcommand{\eqbr}{`=_{\beta}}

\def\coqlibrary#1#2#3{}

% Tutorial

\newboolean{darkfond}
\setboolean{darkfond}{false}
%\setboolean{darkfond}{true}

\definecolor{metavarcolor}{rgb}{0.5,0.0,1.0}
\definecolor{darkgreen}{rgb}{0.1,0.7,0.1}
\definecolor{answercolor}{rgb}{1.0,0.0,0.0}
\definecolor{normalcolor}{rgb}{0.0,0.0,0.0}
\definecolor{exbluecolor}{rgb}{0.1,0.1,0.9}
\definecolor{dontlookcolor}{rgb}{0.5,0.5,0.5}
\definecolor{termcolor}{rgb}{0.0,0.1,0.9}
\definecolor{lookcolor}{rgb}{0.8,0.2,0.0}
\definecolor{prooftermcolor}{rgb}{0.3,0.1,1.0}
\definecolor{typecolor}{rgb}{1.0,0.6,0.0}
\definecolor{taccolor}{rgb}{0.70,0.10,0.0}
\definecolor{pink}{rgb}{0.8,0.6,0.6}
\definecolor{darkmagenta}{rgb}{0.4,0.0,0.6}
\definecolor{darkblue}{rgb}{0.0,0.0,0.6}


\newcommand{\corkscrew}{\boldsymbol{\large \vdash}}
%\newcommand{\disj}{\mbox{$\backslash/$}}
%\newcommand{\conj}{\mbox{$/\backslash$}}
\newcommand{\juge}[3]{\mbox{$#1 \boldsymbol{\vdash} #2 : \color{red}#3 $}}
%\newcommand{\juge}[4]{\mbox{$#1,#2 \corkscrew #3 \boldsymbol{:} #4 $}}
%\newcommand{\smalljuge}[3]{\mbox{$#1 \corkscrew #2 \boldsymbol{:} #3 $}}
%\newcommand{\goal}[3]{\mbox{$#1,#2 \corkscrew^{\!\!\!?} #3  $}}
%\newcommand{\sgoal}[2]{\mbox{$#1\corkscrew^{\!\!\!\!?} #2 $}}
%\newcommand{\reduc}[5]{\mbox{$#1,#2 \corkscrew #3 \rhd_{#4}#5 $}}
%\newcommand{\convert}[5]{\mbox{$#1,#2 \corkscrew #3 =_{#4}#5 $}}
%\newcommand{\convorder}[5]{\mbox{$#1,#2 \corkscrew #3\leq _{#4}#5 $}}
%\newcommand{\wouff}[2]{\mbox{$\emph{WF}(#1)[#2]$}}


\newcommand{\mthese}{\underset{M}{\vdash}}
%\newcommand{\jthese}{\mathop{\vdash}_J}
%\newcommand{\jthese}{\stackrel[J]{}{\vdash}}
\newcommand{\jthese}{\underset{\tiny J}{\vdash}}
%\newcommand{\jthese}{\boldsymbol{\vdash}_{\tiny{ \!\!J}}}

\newcommand{\jsequent}[2]{\mbox{$#1\;\jthese\; #2 $}}
\newcommand{\msequent}[2]{\mbox{$#1\;\mthese\; #2 $}}



%\newcommand{\type}{\boldsymbol{:}}


\newcommand{\answ}{\color{answercolor}}
\newcommand{\coquser}{\color{normalcolor}}
\newcommand{\Ord}{\mbox{$\mathbb{O}$}}
\newcommand{\ords}{\mbox{\(\mathbb{O}\)}}
\newcommand{\nl}{\mbox{$\underset{\scriptstyle NL}{\longrightarrow}$}}
\newcommand{\base}[1]{{\color{lookcolor}#1}}
\newcommand{\om}{\mbox{$\omega$}}

%_{\alpha\beta\delta\iota}
\def\order{\mathcal{O}}
\def\tchecking#1#2#3{\ensuremath{#1 \vdash #2 : #3}}
\def\tconv#1#2#3#4{\ensuremath{#1 \vdash #2 = #3 "~>"
    #4}}
\def\tcumul#1#2#3#4{\ensuremath{#1 \vdash #2 \leq #3 "~>"
    #4}}

\def\tinfer#1#2#3#4#5#6{\ensuremath{#1; #2 \vdash #3 \Uparrow\ 
    "~>" #4 \vdash #5 : #6}}

\def\tcheck#1#2#3#4#5#6#7{\ensuremath{#1; #2 \vdash #3 \Downarrow #4
    "~>" #5 \vdash #6 : #7}}

\usepackage{mathpartir}
%\newcommand{\trule}[1]{\textsc{#1}}
%\newcommand{\irule}[3]
%{\inferrule{ #2 }{ #3 }{\ \trule{#1}}}

\begin{document}


\setbeamertemplate{background canvas}[vertical shading][top=rouge1,middle=rouge1,bottom=rouge1]
\setbeamertemplate{footline}{\hspace{5em} \textcolor{white} {\null \hfill}\hspace{2em}\null \vspace*{3pt}}

\begin{frame}

\begin{textblock*}{40mm}[0,0](10mm,0mm)
 \includegraphics[width=6cm]{INRIA_CHERCHEURS_UK_RVB}
  \end{textblock*}

\begin{textblock*}{8cm}(13mm,50mm)
{\textcolor{white} {
{\huge \thetitle}\\[2mm]
{\theauthor}}}
\end{textblock*}


   \begin{textblock*}{40mm}[0,0](10mm,76mm)
  \begin{picture}(5,80)
\put(0,23){\includegraphics[width=4cm,height=1.5cm]{logobasrougeV1}}
\put(20,50){\tiny \textcolor{rouge2}{Project Team $πr^2$}}
\put(20,40){\tiny \textcolor{rouge2}{INRIA \& PPS, Paris 7}}
\end{picture}
\end{textblock*}


\begin{textblock*}{7cm}(55mm,76mm)
{\textcolor{white}{{\theevent}}}
\end{textblock*}

\vspace*{-4pt}
\end{frame}


%%%%%%%%%%%%%%%%%%%%%%%%%%%

\setbeamertemplate{background canvas}[vertical shading][top=rouge1,bottom=rouge1,middle=rouge1]
\setbeamercolor{toto}{fg=blanc,bg=rouge1}

\setbeamertemplate{footline}
{
\begin{beamercolorbox}[wd=1\paperwidth,ht=15.5pt]{toto}
\raisebox{1ex}
{\includegraphics[width=25mm]{logobastrans}}
  \raisebox{2.5ex}
  {\theauthor{} - \thetitle}\hfill 
  \raisebox{2.5ex}
  {\insertframenumber \hspace{5mm} \null }
\end{beamercolorbox}}

\setbeamertemplate{background canvas}[vertical shading][top=white,middle=white,bottom=white]

\def\bulletfail{\alert{\ding{54}}}
\def\bulletcheck{\ding{52}}

%\setbeamercolor{background}{fg=red,bg=white,text=blue}

%\setbeamercolor{subsection in toc shaded}{fg=gray,bg=white,text=blue}
\setbeamercolor{section in toc}{fg=black,bg=white,text=blue}
\setbeamercolor{subsection in toc}{fg=black,bg=white,text=blue}
%\setbeamertemplate{table of contents shaded}[default]
%\setbeamertemplate{subsection in toc shaded}{\textcolor{gray}}

\setbeamertemplate{section in toc shaded}
{\begin{colormixin}{20!white}{\usebeamertemplate{section in toc}}\end{colormixin}\unskip}

\setbeamertemplate{subsection in toc shaded}
{\begin{colormixin}{20!white}{\usebeamertemplate{subsection in toc}}\end{colormixin}\unskip}

\setbeamertemplate{subsubsection in toc shaded}
{\begin{colormixin}{20!white}{\usebeamertemplate{subsubsection in toc}}\end{colormixin}\unskip}

\def\max{\mathsf{max}}

\section{Introduction}

\begin{frame}
  \frametitle{What are universes?}
  
  Universes are the types of \emph{types}, e.g:
  \begin{itemize}
  \item $\ind{nat}, \ind{bool} : \Type{0}$
  \item $\Type{0} : \Type{1}$
  \item $\ind{list} : \Type{0} "->" \Type{0}$
  \item $∀ α : \Type{0}, \ind{list}~α : \Type{1}$
  \item $∀ n : \ind{nat}, \{ n = 0 \} + \{ n \neq 0 \} : \Type{0}$
  \end{itemize}

\end{frame}

\begin{frame}
  \frametitle{How are they organised?}

  A \emph{hierarchy} of universes $\Type{0} < \Type{1} < \ldots \Type{n}$ used to avoid the $\Type{} : \Type{}$ paradox (system $U^{-}$,
  first version of \name{Martin-Löf} Type Theory (MLTT)).

  \begin{itemize}
  \item Replicates \name{Russell}'s paradox of 
    $\{ x `| x \notin x \}$, the set of all sets etc....
  \item You can think of
    \Type{0} as sets, \Type{1} as classes etc...
  \end{itemize}

\end{frame}
\begin{frame}
  \frametitle{\Coq's theory}

  We call the sort of $t$ the type of the type of $t$, it is necessarily
  a $\Type{i}$ (see e.g. \textsc{Type-prod}).
  
  \begin{figure}
    \begin{mathpar}
      \irule{Type-intro}{\vdash Γ \quad (i \in \mathbb{N})}{Γ \vdash \Type{i} : \Type{i+1}}

      \irule{Type-prod}{Γ \vdash A : \Type{i}\\Γ, x : A \vdash B : \Type{j}}
      {Γ \vdash Π x : A. B : \Type{\max(i,j)}}

    \end{mathpar}
  \end{figure}

  For our purposes, think of Coq's type theory as a stratified,
  predicative System F + primitive inductive/algebraic datatypes.
  The dependent quantification on terms is not relevant here. 
\end{frame}

\begin{frame}
  \frametitle{Typical ambiguity}

  Working with explicit universe indices is cumbersome, annotations
  pervade definitions and proofs.

  \begin{center}
    $"=>"$ Allow \emph{typical ambiguity} (first used by Russell in Principia).
  \end{center}
  
  Idea: write $\Type{}$ to mean any type that works (keeps the system
  consistent).
  
  \begin{itemize}
  \item On paper: let the reader infer levels for universes and check
    consistency. 
  \item On computer: let the computer infer levels and check consistency
    in the background.
  \end{itemize}

  In practice, translation from a source language with anonymous $\Type{}$s
  to a core theory with $\Type{i}$s only.
\end{frame}

\begin{frame}
  \frametitle{Floating universes}

  But in general many $i$'s can work!

  \[\kw{Definition}~\cst{id}~(A : \Type{})~(a : A) := a.\]

  \[\vdash \cst{id} : \Pi (A : \Type{0}),~A \rightarrow A : \Type{1} \]
  \begin{center}or\end{center}
  \[\vdash \cst{id} : \Pi (A : \Type{1}),~A \rightarrow A : \Type{2} \]

  \begin{center}or \ldots\end{center}
  
  \begin{center}
    $"=>"$ Allow \emph{universe variables}.
  \end{center}

\end{frame}

\begin{frame}
  \frametitle{Floating universes and constraints}

  Consistency is now decided by giving an assignation of natural numbers to universe variables, 
  satisfying \emph{constraints}. New judgment $\vdash_{float}$

  \begin{figure}
    \begin{mathpar}
      \irule{Type-intro}{\vdash_{float} Γ \quad (i,j \in \mathbb{L})}{Γ \vdash_{float}
        \Type{i} : \Type{j} "~>" i < j}

      \irule{Type-prod}{Γ \vdash_{float} A : \Type{i}\\Γ, x : A \vdash B : \Type{j}}
      {Γ \vdash_{float} Π x : A. B : \Type{k} "~>" \max(i,j) \leq k}
    \end{mathpar}
  \end{figure}

\end{frame}

\begin{frame}
  \frametitle{Floating universes and constraints \textrm{II}}

  \begin{proposition}[Correctness]
    \textcolor{white}{If $Γ \vdash_{float} t : T "~>" Θ$ and $Θ$ is satisfiable (with
    assignemnt $σ$) then $Γ[σ] \vdash t[σ] : T[σ]$.}
  \end{proposition}  

  Constraints $"=>"$ graph structure based on union-find. $l < l'$
  is an arc, consistency is cycle-freeness, easy to check incrementaly.
\end{frame}

\begin{frame}
  \frametitle{Without polymorphism}

  Floating levels give a false sense of polymorphism:

  \[\kw{Definition}~\cst{id}~(A : \Type{})~(a : A) := a\]

  \[\vdash \cst{id} : \Pi (A : \Type{l}),~A \rightarrow A :
  \Type{\max(l+1,l)} `= \Type{l+1} \]

  $"=>"$ $l$ is \emph{not} quantified at the definition level here, it is \emph{global}: 

  \[\not\vdash \cst{id}~(\Pi (A : \Type{l}),~A \rightarrow A)~\cst{id} : (\Pi (A :
  \Type{l}),~A \rightarrow A)\]

  Because $l+1 \not\le l$.

\end{frame}


\begin{frame}
  \frametitle{With polymorphism}

  Real, bounded polymorphism:

  \[\kw{Polymorphic}~\kw{Definition}~\cst{id}~(A : \Type{})~(a : A) := a\]

  \[\cst{id} : \forall~l, \Pi (A : \Type{l}),~A \rightarrow A\]

  $"=>"$ $l$ is quantified at the definition level now, but second-class, we need
  \emph{instantiation}:

  \[\vdash_{poly} \cstu{id}{k}~(\Pi (A : \Type{l}),~A \rightarrow A)~\cstu{id}{l} : (\Pi (A :
  \Type{l}),~A \rightarrow A) "~>" k < l\]
  
\end{frame}

\section{The current setup}
\frame<beamer>{\frametitle{\thetitle}\tableofcontents[sectionstyle=show/show,subsectionstyle=show/show/show]}


\subsection{Definitions}

\begin{frame}
  \frametitle{Kernel Type Inference}

  Implicit universes with cumulativity, à la Russell.
  
  \emph{In the kernel}, build up a set of universe constraints $Θ$.

  \begin{figure}
  \begin{mathpar}

    \irule{Prod}
    {\tchecking{\Gamma; Θ}{T}{\Type{i}} "~>" Θ_1 \\
     \tchecking{\Gamma, x : T; Θ_1}{U}{\Type{j}} "~>" Θ_2}
    {\tchecking{Γ; Θ}{Π x : T. U}{\Type{\max(i, j)}} "~>" Θ_2}

    \irule{Conv}
    {\tchecking{\Gamma; Θ}{t}{U} "~>" Θ_1 \\
     \tchecking{\Gamma; Θ_1}{V}{s} "~>" Θ_2 \\
     \tcumul{Θ_2}{U}{V}{Θ_3}}
    {\tchecking{Γ; Θ}{t}{V} "~>" Θ_3}
  \end{mathpar}
\end{figure}
\end{frame}

\begin{frame}
  \frametitle{Kernel Conversion}

  \begin{figure}
    \begin{mathpar}
  \irule{Cumul-Sort}
    {}
    {\tcumul{Θ}{\Type{i}}{\Type{j}}{Θ \cup i \leq j}}

    \irule{Cumul-Prod}
    {\tconv{Θ}{U}{U'}{Θ_1} \\
      \tcumul{Θ_1}{T}{T'}{Θ_2}}
    {\tcumul{Θ}{Π x : U. T}{Π x : U'. T'}{Θ_2}}

  \end{mathpar}
\end{figure}  
\end{frame}


\begin{frame}
  \frametitle{Some definitions}

  Algebraic universes and constraints:

  \[\begin{array}{llcl}
    \text{levels} & i, j, le, lt & \in & \mathbb{N} \cup \{ \Prop, \Set \} \\
    \text{universes} & u, v & ::= & i `| \max(\vec{le}, \vec{lt}) \\
    \text{successor} & i+1 & ::= & \max([], i) \\
    \text{order} & \order & ::= & =\ `|\ <\ `|\ \le\ \\

    \text{atomic constraint} & c & ::= & i\ \order\ j \\
    \text{constraints} & Θ & ::= & ε `| c \cup Θ
  \end{array}\]
  
  \pause
  Only handles constraints of the form $u\ \order\ j$ by
  translation to atomic constraints:

  \[\max(i\ j, k) \leq l "<=>" i \leq l \cup j \leq l \cup k < l\]

  Invariant on typing ensures this is the shape of infered constraints
  (Herbelin, TYPES).

\end{frame}
\subsection{Issues}
\begin{frame}
  \frametitle{Issues}

  \begin{itemize}
  \item Constraints are regenerated at each type checking 
  \item Forces the global, unorderly generation of universe variables.
    Any term going out of the kernel must get refreshed universe
    variables because $\max(i,j)$ shouldn't be fed back to it.
    

    % During typing, $\max$ shouldn't appear on the left, i.e.: $Γ
    % \not\vdash \Type{\max(i,j)} : ?$. But it might appear in a type
    % inferred by the kernel, e.g. as a result of typing $Π x : T. U$.
    % So, when taking a term back from the kernel, we must refresh its
    % variables, changing any $\Type{max(i,j)}$ to $\Type{k}$ for some
    % fresh $k$.

  \item To implement universe polymorphism, must hack directly inside the kernel.
    Done for inductive types for now.
  \end{itemize}
\end{frame}  

\section{The new setup}
\frame<beamer>{\frametitle{Universes in \Coq}\tableofcontents[sectionstyle=show/shaded,subsectionstyle=show/show/shaded]}

\begin{frame}
  \frametitle{Constraint checking}
  
  \[\begin{array}{llcl}
    \text{universe context} & Ψ & ::= & \vec{i} "|=" Θ
  \end{array}\]

  Constraints are generated once at \alert{refinement} time (outside the
  kernel):
  
  Inference: {$\tinfer{Γ}{Ψ}{t}{Ψ'}{t'}{T}$}
  
  Checking: {$\tcheck{Γ}{Ψ}{t}{T}{Ψ'}{t'}{T}$}

  \pause
  \begin{figure}
  \begin{mathpar}
    \irule{Check-Type}{$\tcumul{θ}{\Type{i+1}}{T}{θ'}$}
    {$\tcheck{Γ}{us "|=" θ}{\Type{}}{T}{us, i "|=" θ'}{\Type{i}}{T}$}

    \irule{Infer-Cst}
    {(\cst{id} : T) \in Σ}
    {\tinfer{Γ}{Ψ}{\cst{id}}{Ψ}{\cst{id}}{T}}
  \end{mathpar}

  \pause
  \begin{itemize}
  \item The kernel just checks constraints: $Γ; Ψ \vdash t : T$
  \item All universes and constraints that appear in the derivation
    (including conversions) must be in $Ψ$.
  \end{itemize}

\end{figure}

\end{frame}  


\subsection{Universe polymorphic definitions}
\begin{frame}
  \frametitle{Universe Polymorphic definitions}
  
  Now we can introduce universe polymorphism. 

  Suppose a top-level definition $\cst{id} := t : T$.
  \pause
  \begin{enumerate}
  \item {$\tinfer{Γ}{}{T}{Ψ}{T'}{s}$}
    \pause
  \item {$\tcheck{Γ}{Ψ}{t}{T'}{i "|=" θ}{t}{T'}$}
    \pause
  \item Add $\cst{id} : ∀\ i "|=" θ, T' := t$ to the environment.
  \end{enumerate}

  \vspace{2em}

  $"=>"$ Guiding principle: constants are \emph{transparent},
  indistinguishable from their bodies.

\end{frame}

\begin{frame}
  \frametitle{Universe Polymorphic definitions}

  To use $\cst{id}$, we change elaboration of constants to:
  \begin{mathpar}
    \irule{Infer-Cst}
    {(\cst{id} : ∀\ i "|=" θ, T) \in Σ \\
      \vec{i'} : \vec{i} \notin \vec{u}}
    {\tinfer{Γ}{\vec{u} "|=" Θ}{\cst{id}}{\vec{u}, \vec{i'} "|="
        Θ \cup θ[\vec{i'}/\vec{i}]}{\cst{id}_{\vec{i'}}}{T[\vec{i'}/\vec{i}]}}
  \end{mathpar}

   $"=>"$ Constants now carry their universe substitution/instance.

   $"=>"$ Inductives and constructors treated the same way.
   
\end{frame}

\begin{frame}
  \frametitle{Universe Polymorphic definitions: conversion}

  \begin{figure}
    \begin{mathpar}
      \irule{R-δ-l}
      {\cstu{c}{\vec{i}} "->"_\delta t \\
        \tgenconv{R}{ψ}{t~\vec{a}}{u}}
      {\tgenconv{R}{ψ}{\cstu{c}{\vec{i}}~\vec{a}}{u}}
      
      \irule{R-δ-r}
      {\cstu{c}{\vec{i}} "->"_\delta u \\
        \tgenconv{R}{ψ}{t}{u~\vec{a}}}
      {\tgenconv{R}{ψ}{t}{\cstu{c}{\vec{i}}~\vec{a}}}
    \end{mathpar}

    \begin{mathpar}
      \irule{R-FO}
      {\tgenconv{=}{ψ}{\vec{as}}{\vec{bs}} \\
        \alert{\tconstreq{ψ}{\vec{u}}{\vec{v}}}}
      {\tgenconv{R}{ψ}{\cstu{c}{\vec{u}}~\vec{as}}{\cstu{c}{\vec{v}}~\vec{bs}}}
    \end{mathpar}

  \end{figure}
\end{frame}

\subsection{The good, the bad and the ugly}
\begin{frame}
  \frametitle{Advantages to elaboration}

  \begin{itemize}
  \item Reduced trusted code base: checking vs inference.
  \item Reduced polymorphism-specific code (actually
    no, thanks to backward compatibility).
  \item Avoid diffuse use of global gensym $"=>"$ more functional.
  \item User-level control on generated universes and form of
    constraints (simplification, declaration...).
  \item Mixing polymorphic and monomorphic definitions.
  \end{itemize}
\end{frame}
 
\begin{frame}[fragile]
  \frametitle{Design choices}

  Disadvantage (for me and some of you): unification and tactics must become universe-aware.

\begin{verbatim}
Universes.constr_of_global : 
  global_reference -> constr in_universe_context
\end{verbatim}
   
  Unification of $\cst{id}_i$ and $\cst{id}_j$:
  Syntactic equality of $i$ and $j$? Do nothing?

\end{frame}


\begin{frame}
  \frametitle{First-order unification of universes: a first bold idea}

  Due to (notoriously heuristic) first-order unification/conversion of
  constants\ldots we could get too strict universe constraints.
  $\kw{Definition}~U2 := \Type{i}.$
  $\kw{Definition}~U1 : U2 := \Type{j} "~>" j < i$
  $\kw{Definition}~U0 : U1 := \Type{k} "~>" k < j$
  $\kw{Definition}~U02 : U2 := U0 "~>" k < i$

  \[\cst{id}_j~U02~"~"~\cst{id}_i~U0 "~>" i = j\]
  But:
  $\cst{id}_j~U02 "->"^* (U0 "->" U0)$ and
  $\cst{id}_i~U0 "->"^* (U0 "->" U0)$

  \begin{center}
    $"=>"$ Analyse \emph{variance} of universes to relax first-order
    unification. Variance could help minimization too.
  \end{center}

  E.g. for $\cst{id}_i~t "~" \cst{id}_j~u$, $i$ and $j$ do not have to be
  compared.
  \vfill

  \alert{$"=>"$} But it looks like the analysis is hard, not
  compositional and requires full normalisation of the constant bodies.
\end{frame}

\begin{frame}[fragile]
  \frametitle{Pragmatic, heuristic solution}
  
\begin{verbatim}
Evd.fresh_global : ?rigid:rigid -> env -> evar_map -> 
  global_reference -> evar_map * constr

type rigid = 
  | UnivRigid
  | UnivFlexible of bool (* can be algebraic? *)
\end{verbatim}

  \begin{itemize}
  \item Polymorphic constants get elaborated with flexible argument
    levels.
  \item Typical ambiguity (e.g. \Type{}) creates \alert{rigid} variables. 
  \item User-given levels will be rigid
  \end{itemize}
\end{frame}


\begin{frame}
  \frametitle{Unification with universes}

  \begin{figure}

  $\igenconv{R}{ψ}{t}{u}{ψ'}$: unification of $t$ and $u$ under $ψ$.
  \begin{mathpar}
    \irule{Elab-R-δ-left}
    {\cstu{c}{\vec{i}} "->"_\delta t \\
      \igenconv{R}{ψ}{t~\vec{a}}{u}{ψ'}}
    {\igenconv{R}{ψ}{\cstu{c}{\vec{i}}~\vec{a}}{u}{ψ'}}
    
    \irule{Elab-R-δ-right}
    {\cstu{c}{\vec{i}} "->"_\delta u \\
      \igenconv{R}{ψ}{t}{u~\vec{a}}{ψ'}}
    {\igenconv{R}{ψ}{t}{\cstu{c}{\vec{i}}~\vec{a}}{ψ'}}

    \irule{Elab-R-FO}
    {\igenconv{\alert{\textbf{=}}}{ψ}{\vec{as}}{\vec{bs}}{ψ'} \\
      \iconstreq{ψ'}{\vec{u}}{\vec{v}}{ψ''}}
    {\igenconv{\alert{\textbf{R}}}{ψ}{\cstu{c}{\vec{u}}~\vec{as}}{\cstu{c}{\vec{v}}~\vec{bs}}{ψ'}}
  \end{mathpar}
  
  $\iconstreq{ψ}{i}{j}{ψ'}$: unification of
  universe instances.
  \begin{mathpar}
    \irule{Elab-Univ-Eq}
    {\tconstreq{ψ}{i}{j}}
    {\iconstreq{ψ}{i}{j}{ψ}}
    
    \irule{Elab-Univ-Flexible}
    {{i_{\mathsf{f}} `V j_{\mathsf{f}}} \in \vec{u_s} \\
      {\tconsistent{ψ ∧ i = j}}}
    {\iconstreq{(\vec{u_s} \models ψ)}{i}{j}{ψ ∧ i = j}}
  \end{mathpar}
\end{figure}
\end{frame}  

\begin{frame}
  \frametitle{Design choice}

  Universe instances are levels:
  Suppose \[\cst{id} : ∀ i "|=", Π A : \Type{i}, A → A\]    
  \[Γ = A : \Type{i}, P : \cst{fibration}_{i,j} A \vdash Σ_{i j}\ A\ P :
  \Type{\max(i,j)}\]
  
  Levels only, adding constraint if an algebraic would appear:
  \[Γ; \vec{u} "|=" θ \vdash
  \cst{id}\ (Σ\ A\ P) \Uparrow \vec{u}, k "|=" θ\ \cup\ \max(i,j)
  \leq k \vdash \cst{id}_{k} (Σ_{i j}\ A\ P) \ldots\]
  \pause
\end{frame}

\subsection{Minimizing the ugly}
\begin{frame}
  \frametitle{Minimization}

  \begin{center}That's \emph{a lot} of fresh universe variables!!
  \end{center}

  Typical example:
  \[\tinfer{Γ}{Ψ}{\cst{id}\ \constr{true}}{Ψ \cup i "|="
    \Set \le i}{@\cst{id}_{i}\ \ind{bool}\ \constr{true}}{\ind{bool}}\]

  \pause
  We'd want: $@\cst{id}_{\Set}\ \ind{bool}\ \ind{true} : \ind{bool}$,
  no new universe, no additional constraint, just as general.

  \pause  
  \begin{center}
    $"=>"$ Minimization: compute a minimal set of universe variables.
  \end{center}

  See Cardelli's greedy algorithm for $F^{\leq}$ inference, local type
  inference (Pierce \& Turner).

  \begin{itemize}
  \item Requires no other lower constraints on $i$ ($j\ \order\ i$).
  \item \alert{Only} applies to flexible variables. 
  \end{itemize}
\end{frame}


\begin{frame}
  \frametitle{Minimization, results}

  Correctness proof: easy, preservation of local solutions.

  \begin{center}
    Of course this is \alert{\emph{not}} endangering the consistency of \Coq!
  \end{center}  

  \vspace{2em}

  \begin{theorem}[\emph{Conservativity}]
    \textcolor{white}{Unfolding universe 
    polymorphic definitions gives correct typings in the original
    system. Might just not be the most general ones if minimization did
    anything. For inductives, each instantiation is a new copy.}
  \end{theorem}
\end{frame}

\subsection{Dealing with Prop}
\begin{frame}[fragile]
  \frametitle{Dealing with \Prop}

  Let $\cstu{false}{i} : \Type{i+1} \eqdef (Π A : \Type{i}, A : \Type{\max(i+1,i)})$. 

  \vspace{1em}

  But $\cstu{false}{\Prop} "->"^* Π A
  : \Prop, A$, of type \Prop by impredicativity (and \Type{\Prop+1} still). 

  \pause

  \vspace{2em}

  \textbf{Fact:} our universe polymorphism is \emph{incompatible} with the implicit use
  of impredicativity (which has computational content according to
  homotopy models, see hProp's in the HoTT book). The implicit $\Prop{} \le \Type{}$
  rule also causes problems for models and syntax of
  proof-irrelevance (Werner \textit{et al})... 

  \vspace{2em}

  \textbf{Ideal Solution:} Let's get the Rooster and the Syntactic
  Bracket (Herbelin \& Spiwack).

\end{frame}

\begin{frame}[fragile]
  \frametitle{Dealing with \Prop \textrm{II}}

  \textbf{Partial Solutions:}
  \begin{itemize}
  \item Allow to disable $\Prop$ completely, with \verb|-no-prop|, for
    HoTTists.
  \item Disallow instantiating a parameter level $i$ with $\Prop$,
    similar to the current \name{Coq} solution. But get less precise types.
  \item Currently, don't mind the problem, it's ok in practice and you
    don't want to lose the benefit of:

    \[\begin{array}{l}
      \cstu{subset}{i}~(A : \Type{i})~(P : A "->"
      \Prop) : \Type{i} \eqdef \\
      \quad\quad (\Sigma_{i,\Prop}~A~P : \Type{\max(i,\Prop)})
    \end{array}\]
  \end{itemize}
  

\end{frame}

\subsection{Implementation \& benchmarks}
\begin{frame}
  \frametitle{Benchmarks}

  Implementation still in progress but:
  \begin{itemize}
  \item Runs the Homotopy Type Theory Coq library with full universe
    polymophism. No noticeable slowdown. Most definitions
    polymorphic on 6 universes at most.
  \item A universe polymorphic formalization of weak groupoids + an
    interpretation of CC in groupoids, takes 5 min to compile with
    polymorphism and fast projections, impossible without those two
    (inconsistency, exponential slowdown). 1min if deactivating
    universes.
  \end{itemize}
  
\end{frame}

\begin{frame}[fragile]
  \frametitle{Implementation}

  \begin{itemize}
  \item Key technique: hash-consing for fast comparison of universe
    levels, universe instances and algebraic universes. Imperfect as
    deserialization breaks hashconsing (WIP with T. Braibant, PMP, \ldots).

  \item Minimization is fast.
  \item Currently using naive structures, i.e. the kernel side graph
    based on union-find does not do compression and the user-side
    (inference) substitution of universes does not use union-find.
    Optimize last!

  \item Universes have to be normalized at the end of inference now:
    \verb|nf_evars_universes : evar_map -> constr -> constr|

  \end{itemize}
\end{frame}

\section{The past \& the future}
\frame<beamer>{\frametitle{Universes in \Coq}\tableofcontents[sectionstyle=show/shaded,subsectionstyle=show/show/shaded]}

\begin{frame}
  \frametitle{Related work}
  
  \begin{itemize}
  \item Harper and Pollack (TCS'91). Handling of definitions and typical
    ambiguity in type synthesis.
  \item J. Courant: Explicit Universes for CC (TPHOLs'02). 
    User-level declarations of $u \leq i$ in contexts, no other change.
  \item Matita (Coen et al.): checked universes, polymorphism at library
    level.
  \item Pierce and Turner (JFP): Local type inference (based on
    Cardelli's greedy inference algorithm).
  \end{itemize}
  
\end{frame}

\begin{frame}
  \frametitle{Nice things that become possible}

  \begin{itemize}
  \item Universe polymorphic developments: reuse definitions and lemmas
    at different levels.
  \item Polymorphism for universes appearing \emph{inside} structures: 
    old discrepancy between parameters and fields.
  \item Computational relations and rewriting:
    long standing limitation, e.g. for MathClasses.
    Useful for HoTT as well.
  \item Let us \emph{declare} universes and constraints (no user syntax yet).
  \item Resizing rules.
  \end{itemize}
  
\end{frame}



\begin{frame}
  \frametitle{The End}
  \begin{center}
    {\Huge That's all folks!}
  \end{center}
  
\end{frame}


\begin{frame}
  \frametitle{Minimization II}

  At the end of elaboration: $\vec{i} "|=" Θ \vdash t : T$.

  Find a mimimal set of universes variables $\vec{i'} ⊂
  \vec{i}$, universes $\vec{u}$, a substitution $σ : \vec{i}
  "->" \vec{u}$ and constraints $Θ'$ s.t.
  $\vec{i'} "|=" Θ' \cup Θσ$ and $\vec{i'} "|=" Θσ "=>" Θ'$.

  \begin{itemize}
  \item First normalize the constraints w.r.t. loops
    ($l \leq r ∧ r \leq l$) and equalities.
    \pause
  \item Canonicalize $Θ$ w.r.t equalities (except globals)
  \item Mark $i$'s that are fresh universe variables from universe
    instances as candidates for unification + restriction for universes
    ``on the left''.
    \pause
  \end{itemize}
\end{frame}

\begin{frame}
  \frametitle{Minimization III}
  
  We now have $Θ$ with only inequality constraints 
  and a set $f$ of flexible universe variables.
  
  \begin{itemize}
  \item Let $i \in f$, compute its g.l.b: $\max(\vec{j}), j\ \order\ i \in Θ$.
    If $i$ has no lower constraints it must be kept.
  \item Generate upper constraints $\{ glb\ \order\ j `| i\ \order\ j \in Θ \}$
  \item Set $i := glb$ except if $glb$ algebraic and $i$ has upper
    constraints. We can share such $glb$s though.
  \end{itemize}

\end{frame}


\end{document}

%%% Local Variables: 
%%% mode: latex
%%% TeX-PDF-mode: t
%%% TeX-master: "slides"
%%% End: 
